%====================================================
%                PRUEBAS DE CAJA NEGRA
%====================================================

\section{Pruebas de Caja Negra}

%====================================================
%               CU-01  ---> IniciarSesion
%====================================================

\subsection{Caso de prueba: CU01}

En la Tabla \ref{tab:cu01_black_test} se muestra la prueba de caja negra realizada sobre el Caso de Uso: IniciarSesion.

\begin{table}[ht]
\centering
\renewcommand{\arraystretch}{1.3}
\begin{tabular}{|p{3cm}|p{12cm}|}
\hline
\rowcolor[RGB]{191,191,191}
\textbf{Prueba} & Caso de prueba: CU01. \\
\hline
\textbf{Descripción} \cellcolor[RGB]{217,217,217} & Comprobar que un usuario pueda iniciar sesión como propietario y como cliente. \\
\hline
\textbf{Entrada} \cellcolor[RGB]{217,217,217} & Campos del formulario (correo, contraseña). \\
\hline
\textbf{Salida} \cellcolor[RGB]{217,217,217} & El sistema muestra la página principal según el tipo de cuenta. \\
\hline
\textbf{Condición} \cellcolor[RGB]{217,217,217} & El usuario debe tener conexión a internet y haberse registrado previamente. \\
\hline
\end{tabular}
\caption{Caso de prueba: CU01.}
\label{tab:cu01_black_test}
\end{table}

%====================================================
%               CU-02  ---> Registro
%====================================================

\subsection{Caso de prueba: CU02}

\begin{table}[ht]
\centering
\renewcommand{\arraystretch}{1.3}
\begin{tabular}{|p{3cm}|p{12cm}|}
\hline
\rowcolor[RGB]{191,191,191}
\textbf{Prueba} & Caso de prueba: CU02. \\
\hline
\textbf{Descripción} \cellcolor[RGB]{217,217,217} & Comprobar que un usuario pueda crear una cuenta en la aplicación. \\
\hline
\textbf{Entrada} \cellcolor[RGB]{217,217,217} & Datos del formulario de registro (nombre, correo, contraseña). \\
\hline
\textbf{Salida} \cellcolor[RGB]{217,217,217} & El sistema crea la cuenta y redirige a la página de inicio de sesión. \\
\hline
\textbf{Condición} \cellcolor[RGB]{217,217,217} & El correo debe ser válido y no estar registrado previamente. \\
\hline
\end{tabular}
\caption{Caso de prueba: CU02.}
\label{tab:cu02_black_test}
\end{table}

%====================================================
%               CU-03  ---> CerrarSesion
%====================================================

\subsection{Caso de prueba: CU03}

\begin{table}[ht]
\centering
\renewcommand{\arraystretch}{1.3}
\begin{tabular}{|p{3cm}|p{12cm}|}
\hline
\rowcolor[RGB]{191,191,191}
\textbf{Prueba} & Caso de prueba: CU03. \\
\hline
\textbf{Descripción} \cellcolor[RGB]{217,217,217} & Comprobar que el usuario pueda cerrar sesión de su cuenta. \\
\hline
\textbf{Entrada} \cellcolor[RGB]{217,217,217} & Seleccionar opción de cerrar sesión en el menú. \\
\hline
\textbf{Salida} \cellcolor[RGB]{217,217,217} & El usuario es desconectado y redirigido a la pantalla de inicio de sesión. \\
\hline
\textbf{Condición} \cellcolor[RGB]{217,217,217} & El usuario debe estar autenticado en la aplicación. \\
\hline
\end{tabular}
\caption{Caso de prueba: CU03.}
\label{tab:cu03_black_test}
\end{table}

%====================================================
%               CU-04  ---> VerPerfil
%====================================================

\subsection{Caso de prueba: CU04}

\begin{table}[ht]
\centering
\renewcommand{\arraystretch}{1.3}
\begin{tabular}{|p{3cm}|p{12cm}|}
\hline
\rowcolor[RGB]{191,191,191}
\textbf{Prueba} & Caso de prueba: CU04. \\
\hline
\textbf{Descripción} \cellcolor[RGB]{217,217,217} & Comprobar que el usuario pueda visualizar su perfil. \\
\hline
\textbf{Entrada} \cellcolor[RGB]{217,217,217} & Seleccionar la opción de perfil en el menú. \\
\hline
\textbf{Salida} \cellcolor[RGB]{217,217,217} & El sistema muestra los datos del perfil del usuario autenticado. \\
\hline
\textbf{Condición} \cellcolor[RGB]{217,217,217} & El usuario debe estar autenticado en la aplicación. \\
\hline
\end{tabular}
\caption{Caso de prueba: CU04.}
\label{tab:cu04_black_test}
\end{table}

%====================================================
%               CU-05  ---> ModificarPerfil
%====================================================

\subsection{Caso de prueba: CU05}

\begin{table}[ht]
\centering
\renewcommand{\arraystretch}{1.3}
\begin{tabular}{|p{3cm}|p{12cm}|}
\hline
\rowcolor[RGB]{191,191,191}
\textbf{Prueba} & Caso de prueba: CU05. \\
\hline
\textbf{Descripción} \cellcolor[RGB]{217,217,217} & Comprobar que el usuario pueda actualizar sus datos personales. \\
\hline
\textbf{Entrada} \cellcolor[RGB]{217,217,217} & Campos editables del perfil con nuevos datos. \\
\hline
\textbf{Salida} \cellcolor[RGB]{217,217,217} & El sistema actualiza la información del perfil y muestra confirmación. \\
\hline
\textbf{Condición} \cellcolor[RGB]{217,217,217} & El usuario debe estar autenticado y tener datos válidos. \\
\hline
\end{tabular}
\caption{Caso de prueba: CU05.}
\label{tab:cu05_black_test}
\end{table}

%====================================================
%               CU-06  ---> EliminarCuenta
%====================================================

\subsection{Caso de prueba: CU06}

\begin{table}[ht]
\centering
\renewcommand{\arraystretch}{1.3}
\begin{tabular}{|p{3cm}|p{12cm}|}
\hline
\rowcolor[RGB]{191,191,191}
\textbf{Prueba} & Caso de prueba: CU06. \\
\hline
\textbf{Descripción} \cellcolor[RGB]{217,217,217} & Comprobar que el usuario pueda eliminar permanentemente su cuenta. \\
\hline
\textbf{Entrada} \cellcolor[RGB]{217,217,217} & Seleccionar opción de eliminar cuenta y confirmar acción. \\
\hline
\textbf{Salida} \cellcolor[RGB]{217,217,217} & El sistema elimina la cuenta y redirige a la pantalla de inicio de sesión. \\
\hline
\textbf{Condición} \cellcolor[RGB]{217,217,217} & El usuario debe estar autenticado y confirmar la eliminación. \\
\hline
\end{tabular}
\caption{Caso de prueba: CU06.}
\label{tab:cu06_black_test}
\end{table}

%====================================================
%               CU-07  ---> RealizarReserva
%====================================================

\subsection{Caso de prueba: CU07}

\begin{table}[ht]
\centering
\renewcommand{\arraystretch}{1.3}
\begin{tabular}{|p{3cm}|p{12cm}|}
\hline
\rowcolor[RGB]{191,191,191}
\textbf{Prueba} & Caso de prueba: CU07. \\
\hline
\textbf{Descripción} \cellcolor[RGB]{217,217,217} & Comprobar que el usuario pueda realizar una reserva en un negocio. \\
\hline
\textbf{Entrada} \cellcolor[RGB]{217,217,217} & Seleccionar negocio, fecha, hora y confirmar reserva. \\
\hline
\textbf{Salida} \cellcolor[RGB]{217,217,217} & El sistema registra la reserva y muestra confirmación con detalles. \\
\hline
\textbf{Condición} \cellcolor[RGB]{217,217,217} & El negocio debe estar disponible en la fecha y hora seleccionada. \\
\hline
\end{tabular}
\caption{Caso de prueba: CU07.}
\label{tab:cu07_black_test}
\end{table}

%====================================================
%               CU-08  ---> VerReservas
%====================================================

\subsection{Caso de prueba: CU08}

\begin{table}[ht]
\centering
\renewcommand{\arraystretch}{1.3}
\begin{tabular}{|p{3cm}|p{12cm}|}
\hline
\rowcolor[RGB]{191,191,191}
\textbf{Prueba} & Caso de prueba: CU08. \\
\hline
\textbf{Descripción} \cellcolor[RGB]{217,217,217} & Comprobar que el usuario pueda consultar su historial de reservas. \\
\hline
\textbf{Entrada} \cellcolor[RGB]{217,217,217} & Seleccionar sección de reservas pasadas. \\
\hline
\textbf{Salida} \cellcolor[RGB]{217,217,217} & El sistema muestra listado de todas las reservas del usuario. \\
\hline
\textbf{Condición} \cellcolor[RGB]{217,217,217} & El usuario debe tener reservas previas registradas. \\
\hline
\end{tabular}
\caption{Caso de prueba: CU08.}
\label{tab:cu08_black_test}
\end{table}

%====================================================
%               CU-09  ---> CancelarReserva
%====================================================

\subsection{Caso de prueba: CU09}

\begin{table}[ht]
\centering
\renewcommand{\arraystretch}{1.3}
\begin{tabular}{|p{3cm}|p{12cm}|}
\hline
\rowcolor[RGB]{191,191,191}
\textbf{Prueba} & Caso de prueba: CU09. \\
\hline
\textbf{Descripción} \cellcolor[RGB]{217,217,217} & Comprobar que el usuario pueda cancelar una reserva. \\
\hline
\textbf{Entrada} \cellcolor[RGB]{217,217,217} & Seleccionar reserva y confirmar cancelación. \\
\hline
\textbf{Salida} \cellcolor[RGB]{217,217,217} & El sistema cancela la reserva y actualiza el estado. \\
\hline
\textbf{Condición} \cellcolor[RGB]{217,217,217} & La reserva debe estar activa y permitir cancelación. \\
\hline
\end{tabular}
\caption{Caso de prueba: CU09.}
\label{tab:cu09_black_test}
\end{table}

%====================================================
%               CU-10  ---> ExportarEventos
%====================================================

\subsection{Caso de prueba: CU10}

\begin{table}[ht]
\centering
\renewcommand{\arraystretch}{1.3}
\begin{tabular}{|p{3cm}|p{12cm}|}
\hline
\rowcolor[RGB]{191,191,191}
\textbf{Prueba} & Caso de prueba: CU10. \\
\hline
\textbf{Descripción} \cellcolor[RGB]{217,217,217} & Comprobar que el usuario pueda exportar eventos a Calendar. \\
\hline
\textbf{Entrada} \cellcolor[RGB]{217,217,217} & Seleccionar opción de exportar reservas a Google Calendar. \\
\hline
\textbf{Salida} \cellcolor[RGB]{217,217,217} & El sistema sincroniza las reservas con la cuenta de Calendar del usuario. \\
\hline
\textbf{Condición} \cellcolor[RGB]{217,217,217} & El usuario debe tener una cuenta de Google vinculada. \\
\hline
\end{tabular}
\caption{Caso de prueba: CU10.}
\label{tab:cu10_black_test}
\end{table}

%====================================================
%               CU-11  ---> VerNotificaciones
%====================================================

\subsection{Caso de prueba: CU11}

\begin{table}[ht]
\centering
\renewcommand{\arraystretch}{1.3}
\begin{tabular}{|p{3cm}|p{12cm}|}
\hline
\rowcolor[RGB]{191,191,191}
\textbf{Prueba} & Caso de prueba: CU11. \\
\hline
\textbf{Descripción} \cellcolor[RGB]{217,217,217} & Comprobar que el usuario pueda consultar sus notificaciones. \\
\hline
\textbf{Entrada} \cellcolor[RGB]{217,217,217} & Seleccionar sección de notificaciones en la aplicación. \\
\hline
\textbf{Salida} \cellcolor[RGB]{217,217,217} & El sistema muestra el listado de todas las notificaciones recibidas. \\
\hline
\textbf{Condición} \cellcolor[RGB]{217,217,217} & El usuario debe tener notificaciones disponibles. \\
\hline
\end{tabular}
\caption{Caso de prueba: CU11.}
\label{tab:cu11_black_test}
\end{table}

%====================================================
%               CU-12  ---> LeerNotificacion
%====================================================

\subsection{Caso de prueba: CU12}

\begin{table}[ht]
\centering
\renewcommand{\arraystretch}{1.3}
\begin{tabular}{|p{3cm}|p{12cm}|}
\hline
\rowcolor[RGB]{191,191,191}
\textbf{Prueba} & Caso de prueba: CU12. \\
\hline
\textbf{Descripción} \cellcolor[RGB]{217,217,217} & Comprobar que el usuario pueda abrir y marcar como leída una notificación. \\
\hline
\textbf{Entrada} \cellcolor[RGB]{217,217,217} & Seleccionar una notificación de la lista. \\
\hline
\textbf{Salida} \cellcolor[RGB]{217,217,217} & El sistema muestra el contenido y marca como leída. \\
\hline
\textbf{Condición} \cellcolor[RGB]{217,217,217} & Debe existir al menos una notificación disponible. \\
\hline
\end{tabular}
\caption{Caso de prueba: CU12.}
\label{tab:cu12_black_test}
\end{table}

%====================================================
%               CU-13  ---> RecalcularNegocios
%====================================================

\subsection{Caso de prueba: CU13}

\begin{table}[ht]
\centering
\renewcommand{\arraystretch}{1.3}
\begin{tabular}{|p{3cm}|p{12cm}|}
\hline
\rowcolor[RGB]{191,191,191}
\textbf{Prueba} & Caso de prueba: CU13. \\
\hline
\textbf{Descripción} \cellcolor[RGB]{217,217,217} & Comprobar que el sistema muestre negocios según la zona geográfica. \\
\hline
\textbf{Entrada} \cellcolor[RGB]{217,217,217} & Ubicación del usuario o selección manual de zona. \\
\hline
\textbf{Salida} \cellcolor[RGB]{217,217,217} & El sistema muestra los negocios disponibles en esa zona. \\
\hline
\textbf{Condición} \cellcolor[RGB]{217,217,217} & El usuario debe permitir acceso a la ubicación. \\
\hline
\end{tabular}
\caption{Caso de prueba: CU13.}
\label{tab:cu13_black_test}
\end{table}

%====================================================
%               CU-14  ---> AplicarFiltros
%====================================================

\subsection{Caso de prueba: CU14}

\begin{table}[ht]
\centering
\renewcommand{\arraystretch}{1.3}
\begin{tabular}{|p{3cm}|p{12cm}|}
\hline
\rowcolor[RGB]{191,191,191}
\textbf{Prueba} & Caso de prueba: CU14. \\
\hline
\textbf{Descripción} \cellcolor[RGB]{217,217,217} & Comprobar que el usuario pueda aplicar filtros a la búsqueda. \\
\hline
\textbf{Entrada} \cellcolor[RGB]{217,217,217} & Seleccionar criterios de filtro (tipo, horario, calificación). \\
\hline
\textbf{Salida} \cellcolor[RGB]{217,217,217} & El sistema actualiza la lista mostrando solo negocios que cumplan los filtros. \\
\hline
\textbf{Condición} \cellcolor[RGB]{217,217,217} & Deben existir negocios que coincidan con los criterios. \\
\hline
\end{tabular}
\caption{Caso de prueba: CU14.}
\label{tab:cu14_black_test}
\end{table}

%====================================================
%               CU-15  ---> BuscarNegocio
%====================================================

\subsection{Caso de prueba: CU15}

\begin{table}[ht]
\centering
\renewcommand{\arraystretch}{1.3}
\begin{tabular}{|p{3cm}|p{12cm}|}
\hline
\rowcolor[RGB]{191,191,191}
\textbf{Prueba} & Caso de prueba: CU15. \\
\hline
\textbf{Descripción} \cellcolor[RGB]{217,217,217} & Comprobar que el usuario pueda buscar negocios por nombre. \\
\hline
\textbf{Entrada} \cellcolor[RGB]{217,217,217} & Ingresar nombre del negocio en el campo de búsqueda. \\
\hline
\textbf{Salida} \cellcolor[RGB]{217,217,217} & El sistema muestra los negocios que coinciden con la búsqueda. \\
\hline
\textbf{Condición} \cellcolor[RGB]{217,217,217} & Deben existir negocios con ese nombre en la base de datos. \\
\hline
\end{tabular}
\caption{Caso de prueba: CU15.}
\label{tab:cu15_black_test}
\end{table}

%====================================================
%               CU-16  ---> OrientarMapa
%====================================================

\subsection{Caso de prueba: CU16}

\begin{table}[ht]
\centering
\renewcommand{\arraystretch}{1.3}
\begin{tabular}{|p{3cm}|p{12cm}|}
\hline
\rowcolor[RGB]{191,191,191}
\textbf{Prueba} & Caso de prueba: CU16. \\
\hline
\textbf{Descripción} \cellcolor[RGB]{217,217,217} & Comprobar que el sistema oriente el mapa según un punto. \\
\hline
\textbf{Entrada} \cellcolor[RGB]{217,217,217} & Seleccionar un negocio o punto en el mapa. \\
\hline
\textbf{Salida} \cellcolor[RGB]{217,217,217} & El sistema centra y orienta el mapa hacia ese punto. \\
\hline
\textbf{Condición} \cellcolor[RGB]{217,217,217} & El mapa debe estar disponible y visible. \\
\hline
\end{tabular}
\caption{Caso de prueba: CU16.}
\label{tab:cu16_black_test}
\end{table}

%====================================================
%               CU-17  ---> CalcularRuta
%====================================================

\subsection{Caso de prueba: CU17}

\begin{table}[ht]
\centering
\renewcommand{\arraystretch}{1.3}
\begin{tabular}{|p{3cm}|p{12cm}|}
\hline
\rowcolor[RGB]{191,191,191}
\textbf{Prueba} & Caso de prueba: CU17. \\
\hline
\textbf{Descripción} \cellcolor[RGB]{217,217,217} & Comprobar que el sistema calcule la ruta óptima hacia un negocio. \\
\hline
\textbf{Entrada} \cellcolor[RGB]{217,217,217} & Seleccionar opción de ruta para un negocio. \\
\hline
\textbf{Salida} \cellcolor[RGB]{217,217,217} & El sistema muestra la mejor ruta con instrucciones de navegación. \\
\hline
\textbf{Condición} \cellcolor[RGB]{217,217,217} & Debe haber conexión a internet y ubicación disponible. \\
\hline
\end{tabular}
\caption{Caso de prueba: CU17.}
\label{tab:cu17_black_test}
\end{table}

%====================================================
%               CU-18  ---> VerDetallesNegocio
%====================================================

\subsection{Caso de prueba: CU18}

\begin{table}[ht]
\centering
\renewcommand{\arraystretch}{1.3}
\begin{tabular}{|p{3cm}|p{12cm}|}
\hline
\rowcolor[RGB]{191,191,191}
\textbf{Prueba} & Caso de prueba: CU18. \\
\hline
\textbf{Descripción} \cellcolor[RGB]{217,217,217} & Comprobar que el usuario pueda consultar detalles de un negocio. \\
\hline
\textbf{Entrada} \cellcolor[RGB]{217,217,217} & Seleccionar un negocio de la lista. \\
\hline
\textbf{Salida} \cellcolor[RGB]{217,217,217} & El sistema muestra información detallada del negocio. \\
\hline
\textbf{Condición} \cellcolor[RGB]{217,217,217} & El negocio debe existir en la base de datos. \\
\hline
\end{tabular}
\caption{Caso de prueba: CU18.}
\label{tab:cu18_black_test}
\end{table}

%====================================================
%               CU-19  ---> VincularNegocio
%====================================================

\subsection{Caso de prueba: CU19}

\begin{table}[ht]
\centering
\renewcommand{\arraystretch}{1.3}
\begin{tabular}{|p{3cm}|p{12cm}|}
\hline
\rowcolor[RGB]{191,191,191}
\textbf{Prueba} & Caso de prueba: CU19. \\
\hline
\textbf{Descripción} \cellcolor[RGB]{217,217,217} & Comprobar que el propietario pueda vincular un negocio a su cuenta. \\
\hline
\textbf{Entrada} \cellcolor[RGB]{217,217,217} & Datos del negocio a vincular. \\
\hline
\textbf{Salida} \cellcolor[RGB]{217,217,217} & El sistema vincula el negocio y lo asocia a la cuenta del propietario. \\
\hline
\textbf{Condición} \cellcolor[RGB]{217,217,217} & El usuario debe ser propietario y el negocio debe ser válido. \\
\hline
\end{tabular}
\caption{Caso de prueba: CU19.}
\label{tab:cu19_black_test}
\end{table}

%====================================================
%               CU-20  ---> EditarNegocio
%====================================================

\subsection{Caso de prueba: CU20}

\begin{table}[ht]
\centering
\renewcommand{\arraystretch}{1.3}
\begin{tabular}{|p{3cm}|p{12cm}|}
\hline
\rowcolor[RGB]{191,191,191}
\textbf{Prueba} & Caso de prueba: CU20. \\
\hline
\textbf{Descripción} \cellcolor[RGB]{217,217,217} & Comprobar que el propietario pueda modificar su negocio. \\
\hline
\textbf{Entrada} \cellcolor[RGB]{217,217,217} & Campos editables del negocio con nuevos datos. \\
\hline
\textbf{Salida} \cellcolor[RGB]{217,217,217} & El sistema actualiza la información del negocio. \\
\hline
\textbf{Condición} \cellcolor[RGB]{217,217,217} & El propietario debe tener un negocio vinculado. \\
\hline
\end{tabular}
\caption{Caso de prueba: CU20.}
\label{tab:cu20_black_test}
\end{table}

%====================================================
%               CU-21  ---> VerEstadisticas
%====================================================

\subsection{Caso de prueba: CU21}

\begin{table}[ht]
\centering
\renewcommand{\arraystretch}{1.3}
\begin{tabular}{|p{3cm}|p{12cm}|}
\hline
\rowcolor[RGB]{191,191,191}
\textbf{Prueba} & Caso de prueba: CU21. \\
\hline
\textbf{Descripción} \cellcolor[RGB]{217,217,217} & Comprobar que el propietario pueda ver estadísticas de su negocio. \\
\hline
\textbf{Entrada} \cellcolor[RGB]{217,217,217} & Seleccionar sección de estadísticas. \\
\hline
\textbf{Salida} \cellcolor[RGB]{217,217,217} & El sistema muestra métricas y datos del desempeño del negocio. \\
\hline
\textbf{Condición} \cellcolor[RGB]{217,217,217} & El propietario debe tener un negocio con datos disponibles. \\
\hline
\end{tabular}
\caption{Caso de prueba: CU21.}
\label{tab:cu21_black_test}
\end{table}

%====================================================
%               CU-22  ---> MarcarVacaciones
%====================================================

\subsection{Caso de prueba: CU22}

\begin{table}[ht]
\centering
\renewcommand{\arraystretch}{1.3}
\begin{tabular}{|p{3cm}|p{12cm}|}
\hline
\rowcolor[RGB]{191,191,191}
\textbf{Prueba} & Caso de prueba: CU22. \\
\hline
\textbf{Descripción} \cellcolor[RGB]{217,217,217} & Comprobar que el propietario pueda marcar períodos de vacaciones. \\
\hline
\textbf{Entrada} \cellcolor[RGB]{217,217,217} & Seleccionar fechas de inicio y fin de vacaciones. \\
\hline
\textbf{Salida} \cellcolor[RGB]{217,217,217} & El sistema marca esas fechas como no disponibles para reservas. \\
\hline
\textbf{Condición} \cellcolor[RGB]{217,217,217} & Las fechas deben ser válidas y futuras. \\
\hline
\end{tabular}
\caption{Caso de prueba: CU22.}
\label{tab:cu22_black_test}
\end{table}

%====================================================
%               CU-23  ---> DesvincularNegocio
%====================================================

\subsection{Caso de prueba: CU23}

\begin{table}[ht]
\centering
\renewcommand{\arraystretch}{1.3}
\begin{tabular}{|p{3cm}|p{12cm}|}
\hline
\rowcolor[RGB]{191,191,191}
\textbf{Prueba} & Caso de prueba: CU23. \\
\hline
\textbf{Descripción} \cellcolor[RGB]{217,217,217} & Comprobar que el propietario pueda desvincular un negocio. \\
\hline
\textbf{Entrada} \cellcolor[RGB]{217,217,217} & Seleccionar opción de desvincular negocio y confirmar. \\
\hline
\textbf{Salida} \cellcolor[RGB]{217,217,217} & El sistema elimina la asociación del negocio a la cuenta. \\
\hline
\textbf{Condición} \cellcolor[RGB]{217,217,217} & El propietario debe tener un negocio vinculado. \\
\hline
\end{tabular}
\caption{Caso de prueba: CU23.}
\label{tab:cu23_black_test}
\end{table}

%====================================================
%               CU-24  ---> RealizarCitas
%====================================================

\subsection{Caso de prueba: CU24}

\begin{table}[ht]
\centering
\renewcommand{\arraystretch}{1.3}
\begin{tabular}{|p{3cm}|p{12cm}|}
\hline
\rowcolor[RGB]{191,191,191}
\textbf{Prueba} & Caso de prueba: CU24. \\
\hline
\textbf{Descripción} \cellcolor[RGB]{217,217,217} & Comprobar que el propietario pueda agendar una cita manualmente. \\
\hline
\textbf{Entrada} \cellcolor[RGB]{217,217,217} & Datos de la cita (cliente, fecha, hora, servicio). \\
\hline
\textbf{Salida} \cellcolor[RGB]{217,217,217} & El sistema registra la cita en el negocio del propietario. \\
\hline
\textbf{Condición} \cellcolor[RGB]{217,217,217} & El horario debe estar disponible en el negocio. \\
\hline
\end{tabular}
\caption{Caso de prueba: CU24.}
\label{tab:cu24_black_test}
\end{table}

%====================================================
%               CU-25  ---> VerCitas
%====================================================

\subsection{Caso de prueba: CU25}

\begin{table}[ht]
\centering
\renewcommand{\arraystretch}{1.3}
\begin{tabular}{|p{3cm}|p{12cm}|}
\hline
\rowcolor[RGB]{191,191,191}
\textbf{Prueba} & Caso de prueba: CU25. \\
\hline
\textbf{Descripción} \cellcolor[RGB]{217,217,217} & Comprobar que el propietario pueda consultar sus citas programadas. \\
\hline
\textbf{Entrada} \cellcolor[RGB]{217,217,217} & Seleccionar sección de citas. \\
\hline
\textbf{Salida} \cellcolor[RGB]{217,217,217} & El sistema muestra el calendario con todas las citas agendadas. \\
\hline
\textbf{Condición} \cellcolor[RGB]{217,217,217} & El propietario debe tener citas agendadas. \\
\hline
\end{tabular}
\caption{Caso de prueba: CU25.}
\label{tab:cu25_black_test}
\end{table}

%====================================================
%               CU-26  ---> CancelarCita
%====================================================

\subsection{Caso de prueba: CU26}

\begin{table}[ht]
\centering
\renewcommand{\arraystretch}{1.3}
\begin{tabular}{|p{3cm}|p{12cm}|}
\hline
\rowcolor[RGB]{191,191,191}
\textbf{Prueba} & Caso de prueba: CU26. \\
\hline
\textbf{Descripción} \cellcolor[RGB]{217,217,217} & Comprobar que el propietario pueda cancelar una cita. \\
\hline
\textbf{Entrada} \cellcolor[RGB]{217,217,217} & Seleccionar cita y confirmar cancelación. \\
\hline
\textbf{Salida} \cellcolor[RGB]{217,217,217} & El sistema cancela la cita y libera el horario. \\
\hline
\textbf{Condición} \cellcolor[RGB]{217,217,217} & La cita debe estar pendiente y permitir cancelación. \\
\hline
\end{tabular}
\caption{Caso de prueba: CU26.}
\label{tab:cu26_black_test}
\end{table}

%====================================================
%               CU-27  ---> MarcarFavorito
%====================================================

\subsection{Caso de prueba: CU27}

\begin{table}[ht]
\centering
\renewcommand{\arraystretch}{1.3}
\begin{tabular}{|p{3cm}|p{12cm}|}
\hline
\rowcolor[RGB]{191,191,191}
\textbf{Prueba} & Caso de prueba: CU27. \\
\hline
\textbf{Descripción} \cellcolor[RGB]{217,217,217} & Comprobar que el cliente pueda marcar un negocio como favorito. \\
\hline
\textbf{Entrada} \cellcolor[RGB]{217,217,217} & Seleccionar opción de marcar como favorito en un negocio. \\
\hline
\textbf{Salida} \cellcolor[RGB]{217,217,217} & El sistema agrega el negocio a la lista de favoritos del cliente. \\
\hline
\textbf{Condición} \cellcolor[RGB]{217,217,217} & El negocio no debe estar marcado como favorito previamente. \\
\hline
\end{tabular}
\caption{Caso de prueba: CU27.}
\label{tab:cu27_black_test}
\end{table}

%====================================================
%               CU-28  ---> DesmarcarFavorito
%====================================================

\subsection{Caso de prueba: CU28}

\begin{table}[ht]
\centering
\renewcommand{\arraystretch}{1.3}
\begin{tabular}{|p{3cm}|p{12cm}|}
\hline
\rowcolor[RGB]{191,191,191}
\textbf{Prueba} & Caso de prueba: CU28. \\
\hline
\textbf{Descripción} \cellcolor[RGB]{217,217,217} & Comprobar que el cliente pueda desmarcar un negocio de favoritos. \\
\hline
\textbf{Entrada} \cellcolor[RGB]{217,217,217} & Seleccionar opción de desmarcar favorito en un negocio. \\
\hline
\textbf{Salida} \cellcolor[RGB]{217,217,217} & El sistema elimina el negocio de la lista de favoritos. \\
\hline
\textbf{Condición} \cellcolor[RGB]{217,217,217} & El negocio debe estar en la lista de favoritos. \\
\hline
\end{tabular}
\caption{Caso de prueba: CU28.}
\label{tab:cu28_black_test}
\end{table}

\newpage

%====================================================
%           TABLAS DE RESUMEN DE CASOS DE USO
%====================================================

\section{Resumen de Casos de Uso}

\subsection{Casos de uso de Cliente y Propietario}

\begin{table}[h]
\centering
\begin{tabular}{|c|c|c|p{7cm}|}
\hline
\textbf{ID}\cellcolor[RGB]{217,217,217} &
\textbf{Caso de Uso}\cellcolor[RGB]{217,217,217} &
\textbf{Actor}\cellcolor[RGB]{217,217,217} &
\multicolumn{1}{|c|}{\textbf{Descripción}\cellcolor[RGB]{217,217,217}} \\
\hline

CU01 & IniciarSesion & Cliente y Propietario & El usuario inicia sesión en la aplicación. \\
\hline
CU02 & Registro & Cliente y Propietario & El usuario crea una cuenta en la aplicación. \\
\hline
CU03 & CerrarSesion & Cliente y Propietario & El usuario cierra sesión de su cuenta. \\
\hline
CU04 & VerPerfil & Cliente y Propietario & El usuario visualiza su perfil. \\
\hline
CU05 & ModificarPerfil & Cliente y Propietario & El usuario actualiza sus datos personales. \\
\hline
CU06 & EliminarCuenta & Cliente y Propietario & El usuario elimina su cuenta del sistema. \\
\hline
CU07 & RealizarReserva & Cliente y Propietario & El usuario realiza una reserva en un negocio. \\
\hline
CU08 & VerReservas & Cliente y Propietario & El usuario consulta sus reservas pasadas. \\
\hline
CU09 & CancelarReserva & Cliente y Propietario & El usuario cancela una reserva previamente realizada. \\
\hline
CU10 & ExportarEventos & Cliente y Propietario & El usuario exporta eventos a Calendar. \\
\hline
CU11 & VerNotificaciones & Cliente y Propietario & El usuario consulta las notificaciones recibidas. \\
\hline
CU12 & LeerNotificacion & Cliente y Propietario & El usuario abre y marca como leída una notificación. \\
\hline
CU13 & RecalcularNegocios & Cliente y Propietario & El sistema muestra negocios por zona geográfica. \\
\hline
CU14 & AplicarFiltros & Cliente y Propietario & El usuario refina la búsqueda de negocios con filtros. \\
\hline
CU15 & BuscarNegocio & Cliente y Propietario & El usuario busca negocios por nombre. \\
\hline
CU16 & OrientarMapa & Cliente y Propietario & El sistema orienta el mapa hacia un punto. \\
\hline
CU17 & CalcularRuta & Cliente y Propietario & El sistema calcula la mejor ruta hacia un negocio. \\
\hline
CU18 & VerDetallesNegocio & Cliente y Propietario & El usuario consulta información detallada de un negocio. \\
\hline

\end{tabular}
\caption{Casos de uso de Cliente y Propietario}
\label{tab:cliente_propietario}
\end{table}

\newpage

\subsection{Casos de uso de Propietario}

\begin{table}[h]
\centering
\begin{tabular}{|c|c|c|p{6cm}|}
\hline
\textbf{ID}\cellcolor[RGB]{217,217,217} &
\textbf{Caso de Uso}\cellcolor[RGB]{217,217,217} &
\textbf{Actor}\cellcolor[RGB]{217,217,217} & 
\multicolumn{1}{|c|}{\textbf{Descripción}\cellcolor[RGB]{217,217,217}} \\
\hline

CU19 & VincularNegocio & Propietario & El propietario vincula un negocio a su cuenta. \\
\hline
CU20 & EditarNegocio & Propietario & El propietario modifica la información de su negocio. \\
\hline
CU21 & VerEstadisticas & Propietario & El propietario consulta estadísticas de su negocio. \\
\hline
CU22 & MarcarVacaciones & Propietario & El propietario marca períodos de indisponibilidad. \\
\hline
CU23 & DesvincularNegocio & Propietario & El propietario desvincula un negocio de su cuenta. \\
\hline
CU24 & RealizarCitas & Propietario & El propietario agenda una cita manualmente. \\
\hline
CU25 & VerCitas & Propietario & El propietario consulta sus citas programadas. \\
\hline
CU26 & CancelarCita & Propietario & El propietario cancela una cita agendada. \\
\hline

\end{tabular}
\caption{Casos de uso de Propietario}
\label{tab:propietario}
\end{table}

\newpage

\subsection{Casos de uso de Cliente}

\begin{table}[h]
\centering
\begin{tabular}{|c|c|c|p{6cm}|}
\hline
\textbf{ID}\cellcolor[RGB]{217,217,217} &
\textbf{Caso de Uso}\cellcolor[RGB]{217,217,217} &
\textbf{Actor}\cellcolor[RGB]{217,217,217} &
\multicolumn{1}{|c|}{\textbf{Descripción}\cellcolor[RGB]{217,217,217}} \\
\hline

CU27 & MarcarFavorito & Cliente & El cliente marca un negocio como favorito. \\
\hline
CU28 & DesmarcarFavorito & Cliente & El cliente elimina un negocio de sus favoritos. \\
\hline

\end{tabular}
\caption{Casos de uso de Cliente}
\label{tab:cliente}
\end{table}

